\documentclass[11pt,a4paper]{../_shared/altacv}

\geometry{left=1.5cm,right=1.5cm,top=1.cm,bottom=1.2cm,columnsep=1.5cm}

\usepackage[utf8]{inputenc}
\usepackage[T1]{fontenc}
\usepackage[french]{babel}
\usepackage{paracol}
\usepackage{xcolor}
\usepackage{hyperref}
\usepackage{fontawesome5}
\usepackage{setspace}

% Couleurs exactes du CV
\definecolor{SlateGrey}{HTML}{2E2E2E}
\definecolor{LightGrey}{HTML}{666666}
\definecolor{DarkPastelRed}{HTML}{8AC0D9}
\definecolor{PastelRed}{HTML}{00B0DA}
\definecolor{GoldenEarth}{HTML}{B4AD0C}
\colorlet{name}{black}
\colorlet{tagline}{PastelRed}
\colorlet{heading}{DarkPastelRed}
\colorlet{headingrule}{GoldenEarth}
\colorlet{subheading}{PastelRed}
\colorlet{accent}{PastelRed}
\colorlet{emphasis}{SlateGrey}
\colorlet{body}{LightGrey}

\renewcommand{\familydefault}{\sfdefault}

\name{BOILEAU Guillaume}
\personalinfo{
  \email{guillaume.boileaupro@gmail.com}
  \phone{+33 6 59 94 85 84}
  \location{Alpes Maritimes, France}
  \linkedin{guillaume-boileau-phd-891b81265}
  \researchgate{Guillaume-Boileau-2}
  \orcid{0000-0002-3576-6968}
}

\setlength{\parskip}{0.75\baselineskip}

\begin{document}
\makecvheader
\vspace{-0.5cm}
\cvsection{Lettre de motivation}
\vspace{-0.3cm}

{\color{accent}\textbf{Objet :}} \textit{Candidature au poste de Project Manager -- Chef de Projet Logiciels Embarqués Spatial}

{\color{body}

Madame, Monsieur,

Je vous adresse ma candidature au poste de Project Manager Chef de Projet Logiciels Embarqués Spatial au sein d'Alpha-X.

Votre offre retient particulièrement mon attention car elle correspond à un rôle de pilotage technique en environnement spatial exigeant, au croisement de l'ingénierie logicielle, de la coordination d'activités logicielles complexes, du management technique et de la sécurisation de l'exécution. La possibilité d'intervenir sur des projets à fort contenu technologique, impliquant des problématiques d'architecture, d'interfaces et d'entrées système, de structuration des pratiques d'ingénierie, de qualité logicielle, de suivi de l'avancement et de la performance, ainsi que de maîtrise des risques, correspond pleinement au type de responsabilités que je souhaite aujourd'hui exercer.

Mon expérience m'a conduit à évoluer sur des projets scientifiques et techniques de grande ampleur, dans lesquels la gestion de projet, le pilotage d'équipes techniques, la coordination transverse, la structuration des activités, la maîtrise de la qualité et la sécurisation de l'exécution occupaient une place centrale, en cohérence directe avec les responsabilités attendues pour le poste de Chef de Projet Logiciels Embarqués Spatial. Dans le cadre de LISA, future grande mission spatiale de l'Agence spatiale européenne consacrée à la détection des ondes gravitationnelles dans l'espace, la France porte un engagement structurant sur le segment sol scientifique, en particulier pour la production, la comparaison et la consolidation des catalogues. Dans ce contexte, l'enjeu est de transformer des données instrumentales complexes et les sorties probabilistes de plusieurs pipelines en catalogues scientifiques robustes, traçables et exploitables. J'y ai assuré un rôle concret de gestion de projet et de gestion d'équipe, en organisant les activités associées, en structurant le travail collectif, en coordonnant les contributions, en encadrant des stagiaires, en pilotant des équipes d'ingénieurs impliquées dans le développement, et en veillant à la cohérence entre besoins fonctionnels, interfaces, contraintes de développement et exigences de qualité. Avec CATpyLISA, dont j'ai conçu l'architecture logique et réalisé les principaux développements, j'ai structuré un cadre permettant de comparer, valider et consolider plusieurs sorties L2 issues de global fits afin de produire des catalogues L3 traçables, robustes et scientifiquement exploitables. Retenu comme base du pipeline L3, ce projet constitue un maillon critique du segment sol scientifique de LISA et participe directement à l'un des engagements français structurants dans cette mission européenne ; il m'a conduit à faire converger des besoins techniques multiples, à suivre l'avancement des travaux, à mettre sous contrôle la qualité et la robustesse des développements, à maintenir la cohérence d'ensemble entre objectifs scientifiques, contraintes de développement, exigences de qualité et résultats attendus, et à porter un projet logiciel au croisement de l'architecture, de la coordination technique et de la montée en maturité collective.

Le projet CATpyLISA illustre particulièrement bien cette dimension de gestion de projet. J'en suis à l'origine et j'en ai conçu l'architecture logique ainsi que les principaux développements. Ce projet a été structuré pour répondre à un besoin critique de la mission LISA : comparer, valider et consolider plusieurs sorties L2 issues de global fits afin de produire des catalogues L3 traçables, robustes et scientifiquement exploitables. Au-delà du développement logiciel lui-même, cela m'a conduit à porter un cadre de travail collectif, à formaliser un modèle de données commun, à structurer les méthodes de comparaison entre pipelines, à définir des critères d'évaluation et des indicateurs partagés, et à organiser la montée en maturité et l'industrialisation progressive du projet dans un cadre collaboratif. Ce travail, retenu comme base du pipeline L3, m'a placé dans une position de pilotage technique, à l'interface entre besoins métier, architecture logicielle, coordination des contributeurs et exigences de qualité.

Cette expérience m'a permis de développer une approche très concrète du management de projet technique. J'ai l'habitude de clarifier les objectifs, découper les activités, organiser les priorités, suivre l'avancement, identifier les points de blocage, challenger les interfaces et sécuriser les livrables. J'accorde une grande importance à la robustesse des workflows, à la traçabilité des résultats, à la cohérence entre briques logicielles, au suivi de la qualité et de la performance, ainsi qu'à la capacité d'une équipe à progresser dans un cadre structuré. J'ai également l'habitude de travailler avec plusieurs interlocuteurs techniques, de maintenir l'alignement entre contributions, d'accompagner la progression collective et de faire évoluer un projet dans un environnement où les exigences méthodologiques et qualitatives sont fortes.

Mon parcours m'a aussi donné une forte sensibilité aux sujets de fiabilité, de risques techniques, de dette méthodologique et de performance globale des chaînes de traitement. Dans les projets sur lesquels j'ai travaillé, il ne s'agissait pas seulement de produire du code, mais de garantir que les résultats soient comparables, exploitables, défendables et maintenables dans la durée. Cette culture de la qualité, du suivi et de la sécurisation de l'exécution constitue l'un des aspects centraux de ma façon de piloter un projet.

Plus récemment, mon expérience en développement logiciel pour des infrastructures de recharge de véhicules électriques m'a permis de renforcer ma compréhension des problématiques d'intégration, d'interfaces système, de flux opérationnels, de supervision et de déploiement dans un cadre plus directement industriel. Cette double expérience, dans le spatial scientifique et dans un environnement logiciel plus applicatif, me permet d'aborder les projets avec une vision à la fois structurée, technique et orientée exécution.

Je suis particulièrement attiré par Alpha-X pour son positionnement au plus près des enjeux réels du spatial, sa culture de l'expertise et du pilotage, ainsi que par la nature même du poste, qui demande d'organiser les activités logicielles, d'accompagner les équipes, de structurer les pratiques d'ingénierie, de maîtriser les interfaces, de suivre la qualité et la performance, et de piloter les risques. Je me reconnais pleinement dans cette logique de responsabilité, de coordination et de progression collective.

Je souhaite aujourd'hui mettre à profit mon expérience en pilotage technique, en structuration de projet logiciel complexe et en environnement spatial pour contribuer à des projets nécessitant méthode, exigence, autonomie et leadership. Je serais très heureux de pouvoir échanger avec vous afin de vous présenter plus en détail mon parcours, ma motivation et la manière dont je pourrais contribuer aux projets portés par Alpha-X.

Veuillez agréer, Madame, Monsieur, l'expression de mes salutations distinguées.

\textbf{Guillaume Boileau}

}

\end{document}